\documentclass[12pt,a4paper]{amsart}

\usepackage[a4paper,textwidth=160mm,textheight=215mm]{geometry}
\usepackage{amsfonts,amssymb,amsthm,amsmath,mathrsfs}
\PassOptionsToPackage{no-math}{fontspec}
\usepackage{xeCJK}
\usepackage[hidelinks]{hyperref}

\setCJKmainfont[BoldFont={SimHei},ItalicFont={KaiTi}]{SimSun}
\setCJKsansfont{SimHei}
\setCJKmonofont{SimSun}

\linespread{1.30}
\setlength{\parindent}{2em}
\numberwithin{equation}{section}

\theoremstyle{plain}
\newtheorem{theorem}{定理}[section]
\newtheorem{proposition}[theorem]{命题}
\newtheorem{lemma}[theorem]{引理}
\newtheorem{corollary}[theorem]{推论}
\theoremstyle{definition}
\newtheorem{remark}[theorem]{注}
\newtheorem{definition}[theorem]{定义}

\renewcommand{\abstractname}{摘要}
\renewcommand{\keywordsname}{关键词}
\renewcommand{\subjclassname}{\textup{2020} 年数学主题分类}
\renewcommand{\proofname}{证明}
\renewcommand{\refname}{参考文献}

\newcommand{\HS}{\mathrm{HS}}
\newcommand{\tr}{\operatorname{tr}}
\newcommand{\dd}{\,\mathrm{d}}
\newcommand{\eps}{\varepsilon}

\begin{document}

\title[对角算子的谱范数]{条件哥德巴赫方法中\\ 对角算子的谱范数}

\author{}

\date{}

\subjclass[2020]{主分类 11N13；次分类 47B35, 47A75}

\keywords{哥德巴赫猜想，对角算子，哈代空间，Schur 检验，加权素数定理}

\begin{abstract}
在 Li 与 Liu 的条件哥德巴赫方法中，加权哥德巴赫误差被精确表示为一个作用在
哈代空间上的算子 $T_N=\mathcal P_N-\mathcal D_N-\mathcal K_N$ 的二次型。本文
确定对角分量 $\mathcal D_N$ 的首阶范数：
\[
 \|\mathcal D_N\|=\Bigl(\frac{\pi}{e}+o(1)\Bigr)\log N .
\]
证明用到 $\mathcal D_N$ 在 Fourier--Laplace 坐标下的精确核、加权素数定理、
一个光滑模型核（其范数由 Schur 检验与匹配的试验向量精确到首阶），以及把模型
范数搬回 $\mathcal D_N$ 的三区域比较。结论是无条件的，并且改进了 \cite{LL}
中对角分量所用的希尔伯特--施密特界。
\end{abstract}

\maketitle

\section{引言}

\subsection*{1.1.\ 对角算子}

在 \cite{LL} 中，Li 与 Liu 证明了二元哥德巴赫猜想的一个条件形式。他们的论证
把加权哥德巴赫误差化归为竖直条带的哈代空间上一个紧自伴算子 $T_N$ 的精确二次型，
且 $T_N$ 分解为
\begin{equation}\label{eq:split}
 T_N=\mathcal P_N-\mathcal D_N-\mathcal K_N,
\end{equation}
其中 $\mathcal D_N$ 是精确的对角算子，$\mathcal K_N$ 是贝塔模型，而
$\mathcal P_N$ 是\emph{奇特征相关算子}（它是形如 $UU^*$ 的正算子，但不是正交投影）。
本文确定对角分量的首阶范数。

设 $N>e^{10}$ 为偶数。我们沿用 \cite[\S5--6]{LL} 的记号：
\[
 \alpha_N=\frac1{\log N},\qquad \beta=\frac18,\qquad
 S_N=\{s:\alpha_N<\Re s<\beta\},\qquad
 \Gamma_N=(\alpha_N+i\mathbb R)\cup(\beta+i\mathbb R),
\]
而 $\mathcal H_N=H^2(S_N)$ 是条带的哈代空间，带边界范数
\begin{equation}\label{eq:inner}
 \langle h,g\rangle_{\mathcal H_N}
 =\int_{\Gamma_N}h(s)\overline{g(s)}\,|\mathrm ds|;
\end{equation}
由 \cite[引理 5.2]{LL}，边界值在 $L^2$ 中存在，且 \eqref{eq:inner} 是
$\mathcal H_N$ 的内积。内积对第一个变量线性，我们记
$(f\otimes g)h=\langle h,g\rangle f$。

对角算子是正秩一级数
\begin{equation}\label{eq:DN}
 \mathcal D_N=\sum_{(n,N)=1}\frac{\Lambda^2(n)}{n}\,d_n\otimes d_n,
 \qquad d_n(s)=\frac{n^{-s}}{1+s},
\end{equation}
其中 $\Lambda$ 是冯·曼戈尔特函数。等价地，记 $q$ 为与 $N$ 互素的素数幂，
\begin{equation}\label{eq:phiq}
 \varphi_q(s)=\frac{q^{-1-s}}{1+s}=q^{-1}d_q(s),
\end{equation}
则
\begin{equation}\label{eq:DN2}
 \mathcal D_N=\sum_{q}q\Lambda(q)^2\,\varphi_q\otimes\varphi_q .
\end{equation}
该级数在迹范数下收敛，且 $\operatorname{tr}\mathcal D_N\ll\alpha_N^{-2}$；
见 \cite[命题 6.5]{LL}。该命题只用到迹界，对算子范数而言它给出
\begin{equation}\label{eq:coarse}
 \|\mathcal D_N\|\le\|\mathcal D_N\|_{\HS}\le\operatorname{tr}\mathcal D_N
 \ll(\log N)^2 .
\end{equation}

\subsection*{1.2.\ 结果陈述}

\begin{theorem}\label{thm:main}
设 $N>e^{10}$ 为偶数，$\mathcal D_N$ 为对角算子 \eqref{eq:DN}。则
$\mathcal D_N$ 是自伴、正、迹类的，且
\begin{equation}\label{eq:main}
 \|\mathcal D_N\|=\Bigl(\frac{\pi}{e}+o(1)\Bigr)\log N .
\end{equation}
\end{theorem}

\begin{corollary}\label{cor:HS}
有 $\|\mathcal D_N\|_{\HS}\ll(\log N)^{3/2}$。
\end{corollary}

\eqref{eq:main} 中的算子范数比迹界 \eqref{eq:coarse} 小一个 $\log N$ 量级的
因子，比推论~\ref{cor:HS} 的希尔伯特--施密特界小一个 $(\log N)^{1/2}$ 因子。
常数 $\pi/e$ 不是方法的产物：它是 $\mathcal D_N$ 在主谱方向上的局部极限——
一个卷积核——的范数。

\begin{remark}\label{rem:eigenvector}
\emph{（实际被证明的是什么。）} 定理~\ref{thm:main} 的证明给出一族\emph{显式的
局域化近极大向量}，而不是顶特征向量。具体地，记
\begin{equation}\label{eq:peakpoint}
 y_D^*=\frac12+\frac1{2\alpha_N},
\end{equation}
并取 $B=\alpha_N^{-1/2}$，则归一化的区间试验函数
$f=B^{-1/2}\mathbf 1_{[y_D^*-B/2,\,y_D^*+B/2]}$ 满足
\[
 \langle\mathcal D_Nf,f\rangle\ge\Bigl(\frac{\pi}{e}-o(1)\Bigr)\alpha_N^{-1};
\]
这正是第~\ref{sec:model} 节的匹配下界与定理~\ref{thm:main} 合起来给出的结论。
于是算子范数在相差 $o(\alpha_N^{-1})$ 的意义下由这样的向量取到：它们的
Fourier--Laplace 质量局域在 $y_D^*$ 附近、宽度 $\ll\alpha_N^{-1/2}$，即落在
$\log q\approx\frac12\log N$ 的素数幂上。

必须强调仍然存在的逻辑缺口。要\emph{识别}真正的顶特征向量（以及它的高斯轮廓），
还需两项输入：(i) 顶端的谱隙
$\lambda_0(\mathcal D_N)-\lambda_1(\mathcal D_N)\gg1$；(ii) 一个谱投影
（或 Mosco 型）论证，说明上述近极大向量收敛到顶谱子空间。本文两者均未证明；
且上面的试验函数是阶梯函数，而特征函数候选是光滑的（在缩放变量
$X=\sqrt{\alpha_N}(y-y_D^*)$ 下形如 Hermite 轮廓 $Ce^{-X^2/\sqrt2}$）。
因此本文的结论只涉及算子范数。
\end{remark}

\begin{remark}\label{rem:numerical}
对下文模型核 \eqref{eq:model} 作离散化，当
$\alpha_N=0.1$、$0.05$、$0.025$、$0.0125$ 时，最大特征值分别为
$\pi/(e\alpha_N)$ 的 $0.730$、$0.850$、$0.923$、$0.972$ 倍；该比值单调
趋于 $1$，与 \eqref{eq:main} 相符。
所用的离散化如下：把半直线 $[y_0,\infty)$ 截断于 $y_0+40/\alpha_N$，等分为
$N_{\rm grid}=4000$ 个子区间；在每个子区间上用中点值代替核，所得的实对称
矩阵用标准 $QR$ 算法对角化；表中数字是最大特征值除以 $\pi/(e\alpha_N)$。
截断误差可忽略，因为核按 $e^{-(1-\alpha_N)|y-u|}$ 衰减，而权重
$e^{-2\alpha_Ny}$ 使 $y-y_0\gtrsim10/\alpha_N$ 之后贡献极小。这些数值只作
说明，不用于证明。
\end{remark}

\subsection*{1.3.\ 方法}

证明分五步。

\begin{enumerate}
\item[\rm(A)] \emph{精确核。} 在 Fourier--Laplace 坐标 $y=\log q$ 下，
$\mathcal D_N$ 成为积分算子，其核为
\[
 \mathscr D_N(y,u)=2\pi W_N(\min(y,u))\sqrt{m(y)m(u)},\qquad
 W_N(y)=\sum_{\substack{q\le e^y\\(q,N)=1}}q\Lambda(q)^2 ,
\]
其中 $m(y)=e^{-2(1+\alpha_N)y}+e^{-2(1+\beta)y}$（第~\ref{sec:kernel} 节）。
\item[\rm(B)] \emph{加权素数定理。} 权重 $W_N$ 在
$y\in[c/\alpha_N,C/\alpha_N]$ 上一致满足
$W_N(y)=\frac12e^{2y}(y-\frac12)(1+o(1))$（第~\ref{sec:pnt} 节）。
\item[\rm(C)] \emph{模型核。} 在该区域中把 $\mathscr D_N$ 换成显式核
\[
 \mathscr A_N(y,u)=\pi\Bigl(\min(y,u)-\frac12\Bigr)
 e^{-2\alpha_N\min(y,u)}e^{-(1+\alpha_N)|y-u|},
\]
其范数由 Schur 检验与匹配的试验向量算出（第~\ref{sec:model} 节）。
\item[\rm(D)] \emph{比较。} 对 $y$ 轴作三区域分解，表明模型范数与
$\mathcal D_N$ 的范数相差 $o(\alpha_N^{-1})$（第~\ref{sec:compare} 节）。
\item[\rm(E)] \emph{合拢。} 由 $\|\mathscr A_N\|\sim\pi/(e\alpha_N)$ 即得定理
（第~\ref{sec:proof} 节）。
\end{enumerate}

\subsection*{1.4.\ 约定}

记 $f\ll g$ 与 $f=O(g)$ 表示 $|f|\le Cg$，$C$ 为绝对常数；下标表示常数依赖。
除非另行说明，所有常数都是绝对的。我们不使用任何未证假设：算子 $\mathcal D_N$
的来源是算术的，但它的核估计是解析的。

\subsection*{1.5.\ 关于本版本的说明}

本版本相对初稿作了如下订正。它们都不改变定理~\ref{thm:main} 的陈述与常数，
但其中两条补上了实质性的论述缺口。

\begin{enumerate}
\item[\rm(i)] 关于顶特征向量的陈述已收窄为证明实际给出的内容：\emph{局域化近极大
向量}的存在性；并明确指出识别顶特征向量还缺两项输入（谱隙与谱投影）。见注~\ref{rem:eigenvector}。
\item[\rm(ii)] \S1.3 中 $W_N$ 的定义已与 \eqref{eq:Dkernel} 统一：求和现在带有
限制 $(q,N)=1$。
\item[\rm(iii)] 在 \S\ref{sec:compare} 中反复使用的估计 $y_0=O(\log\log N)$
现在给出证明而非默认；见 \eqref{eq:y0small} 及其后的几行。
\item[\rm(iv)] 中心块中从核不等式
$|\mathscr D_N-\mathscr A_N|\le\epsilon_N\mathscr A_N$ 到算子不等式的过渡，
现在通过二次型与 $\mathscr A_N$ 的正性写出。见引理~\ref{lem:center} 的证明。
\item[\rm(v)] 关于 $p\mid N$ 贡献的估计，其适用的 $x$ 范围现在在
注~\ref{rem:primepowers} 中显式写出。
\item[\rm(vi)] \S5.7 的交叉块估计现在对 $(-,0)$、$(0,+)$、$(-,+)$ 三种情形分别列出；
此前是合并写的。
\item[\rm(vii)] 数值注现在记录了复现所需的离散化细节（截断点、网格、求积规则、
特征值算法）；见注~\ref{rem:numerical}。
\item[\rm(viii)] 算子 $\mathcal P_N$ 不再称为"精确的特征投影"（它不是正交投影），
而称为\emph{奇特征相关算子}；见 \S1.1。
\end{enumerate}

\section{哈代空间与精确核}\label{sec:kernel}

\subsection*{2.1.\ 哈代空间}

我们回顾 \cite[\S5]{LL} 中 $\mathcal H_N$ 的描述。函数 $h$ 属于 $H^2(S_N)$，
如果它在 $S_N$ 中全纯且
\[
 \|h\|^2_{H^2(S_N)}
 =\sup_{0<\eps<(\beta-\alpha_N)/2}
 \int_{\mathbb R}\bigl(|h(\alpha_N+\eps+it)|^2+|h(\beta-\eps+it)|^2\bigr)\dd t<\infty .
\]
边界值在 $L^2$ 中存在，且 \eqref{eq:inner} 成立。Fourier--Laplace 表示为
\begin{equation}\label{eq:PW}
 h(s)=\frac1{\sqrt{2\pi}}\int_{\mathbb R}\psi(\tau)e^{s\tau}\dd\tau,
 \qquad
 \|h\|^2_{\mathcal H_N}=\int_{\mathbb R}|\psi(\tau)|^2w_H(\tau)\dd\tau,
\end{equation}
其中 $w_H(\tau)=e^{2\alpha_N\tau}+e^{2\beta\tau}$（见 \cite[引理 5.2]{LL}）。

\subsection*{2.2.\ 核}

\begin{lemma}\label{lem:kernel}
对 $s,z\in\Gamma_N$，
\begin{equation}\label{eq:kernel}
 \kappa_N(s,z):=\sum_{(n,N)=1}\frac{\Lambda^2(n)}{n}\,d_n(s)\overline{d_n(z)}
 =\sum_q q\Lambda(q)^2\varphi_q(s)\overline{\varphi_q(z)} .
\end{equation}
因此 $\mathcal D_N$ 自伴、正，且
\begin{equation}\label{eq:trace}
 \operatorname{tr}\mathcal D_N\ll\alpha_N^{-2}.
\end{equation}
\end{lemma}

\begin{proof}
两个表达式相等，因为 $d_n=n\varphi_n$ 而权重为 $q\Lambda(q)^2$；更精确地，
$q\Lambda(q)^2\varphi_q\otimes\varphi_q=\Lambda(q)^2q^{-1}d_q\otimes d_q$。
该算子是系数非负的秩一算子 $\varphi\otimes\varphi$ 之和，故正且自伴。最后，
\[
 \|\varphi_q\|^2=\sum_{\rho\in\{\alpha_N,\beta\}}q^{-2(1+\rho)}
 \int_{\mathbb R}\frac{\dd t}{(1+\rho)^2+t^2}
 =\pi\sum_{\rho\in\{\alpha_N,\beta\}}\frac{q^{-2(1+\rho)}}{1+\rho},
\]
故
\[
 \operatorname{tr}\mathcal D_N
 =\pi\sum_{\rho\in\{\alpha_N,\beta\}}\frac1{1+\rho}
 \sum_q\Lambda(q)^2q^{-1-2\rho}\ll\alpha_N^{-2},
\]
因为当 $0<\rho\le\beta$ 时 $\sum_q\Lambda(q)^2q^{-1-2\rho}\asymp(2\rho)^{-2}$。
\end{proof}

\subsection*{2.3.\ Fourier--Laplace 坐标}

置
\begin{equation}\label{eq:m}
 y=-\tau,\qquad m(y)=e^{-2(1+\alpha_N)y}+e^{-2(1+\beta)y},
\end{equation}
对 $h\in\mathcal H_N$（变换为 $\psi$）令 $f(y)=\psi(-y)\sqrt{w_H(-y)}$。则
$y\mapsto f$ 是从 $\mathcal H_N$ 到 $L^2(\mathbb R,\dd y)$ 的酉映射，且
\begin{equation}\label{eq:unit}
 \|h\|_{\mathcal H_N}=\|f\|_{L^2(\dd y)} .
\end{equation}
设 $y_0=\log q_1$，其中 $q_1$ 是与 $N$ 互素的最小素数幂。由于 $N$ 为偶数，
$q_1\ge3$，故在下文核的支集上
\begin{equation}\label{eq:y0}
 y\ge y_0\ge\log3>\frac12 .
\end{equation}
下面还会多次用到 $y_0$ 相对于 $\log N$ 很小这一事实：
\begin{equation}\label{eq:y0small}
 y_0\le\log\log N+O(1),\qquad\text{即}\quad y_0=O(\log\log N).
\end{equation}
事实上，每个小于 $q_1$ 的素数 $p$ 都整除 $N$（否则 $p$ 本身就是与 $N$ 互素
且小于 $q_1$ 的素数幂），因此 $\prod_{p<q_1}p\mid N$，从而
$\prod_{p<q_1}p\le N$。由素数定理 $\prod_{p<q_1}p=e^{q_1(1+o(1))}$，故
$q_1\le(\log N)(1+o(1))$，即得 \eqref{eq:y0small}。特别地，对充分大的 $N$ 有
$y_0\le\frac12\log N$，故 $y_0$ 总是远低于峰值
$y_D^*=\frac12+\frac1{2\alpha_N}$。

\begin{lemma}\label{lem:FL}
在酉映射 \eqref{eq:unit} 下，算子 $\mathcal D_N$ 成为 $L^2(\dd y)$ 上的积分
算子，核为
\begin{equation}\label{eq:Dkernel}
 \mathscr D_N(y,u)=2\pi\,W_N(\min(y,u))\,\sqrt{m(y)m(u)},
 \qquad
 W_N(y)=\sum_{\substack{q\le e^y\\(q,N)=1}}q\Lambda(q)^2 .
\end{equation}
特别地，$\mathscr D_N(y,u)=0$ 除非 $y,u\ge y_0$。
\end{lemma}

\begin{proof}
对 $s\in\Gamma_N$，
\[
 \varphi_q(s)=\frac{q^{-1-s}}{1+s}=\int_{\log q}^{\infty}e^{-(1+s)v}\dd v
 =\int_{-\infty}^{-\log q}e^{(1+s)\tau}\dd\tau,
\]
故在 \eqref{eq:PW} 中 $\psi_q(\tau)=\sqrt{2\pi}e^{\tau}\mathbf 1_{\tau\le-\log q}$。
在加权表示 $L^2(w_H\dd\tau)$ 中，$\varphi_q\otimes\varphi_q$ 的核为
$2\pi e^{\tau+\tau'}\mathbf 1_{\tau,\tau'\le-\log q}$。经 $\psi\mapsto\psi\sqrt{w_H}$
单位化并令 $y=-\tau$，得
\[
 2\pi e^{-y-u}\sqrt{w_H(-y)w_H(-u)}\mathbf 1_{y,u\ge\log q}
 =2\pi\sqrt{m(y)m(u)}\mathbf 1_{y,u\ge\log q} .
\]
对 $q\le\min(e^y,e^u)$ 以权重 $q\Lambda(q)^2$ 求和即得 \eqref{eq:Dkernel}。
\end{proof}

\section{加权素数定理与模型核}\label{sec:pnt}

\subsection*{3.1.\ 加权素数定理}

\begin{lemma}\label{lem:pnt}
对 $x\ge2$，
\begin{equation}\label{eq:pnt}
 \sum_{p\le x}p(\log p)^2
 =\frac12x^2\log x-\frac14x^2+O\bigl(x^2(\log x)^Ce^{-c\sqrt{\log x}}\bigr),
\end{equation}
其中 $C,c>0$ 为绝对常数。
\end{lemma}

\begin{proof}
由形如
$\theta(x)=\sum_{p\le x}\log p=x+O(xe^{-c\sqrt{\log x}})$ 的素数定理与部分求和，
\[
 \sum_{p\le x}p(\log p)^2=\int_2^x t\log t\,\dd\theta(t)
 =\int_2^x t\log t\,\dd t+O\bigl(x^2(\log x)^Ce^{-c\sqrt{\log x}}\bigr),
\]
而 $\int_2^x t\log t\,\dd t=\frac12x^2\log x-\frac14x^2+O(1)$。
\end{proof}

\begin{remark}\label{rem:primepowers}
高次素数幂 $q=p^a$（$a\ge2$）对 $\sum_{q\le x}q\Lambda(q)^2$ 的贡献为
$\ll x^{3/2}\log^2x$；而限制 $(q,N)=1$ 只去掉整除 $N$ 的素数的幂，其总贡献为
$\ll x\log^2N$。这两个界对 $2\le x\le N^{A}$（$A$ 固定）一致成立——特别地，
在下文唯一用到的范围 $x=e^y$、$y\asymp\log N$ 内成立。（对更大的 $x$，
$\ll x\log^2N$ 已不再有用，但下文从不在那里使用它。）在下文的范围内两者都是
低阶项。
\end{remark}

\subsection*{3.2.\ $W_N$ 的渐近}

\begin{lemma}\label{lem:W}
对固定的 $0<c<C<\infty$ 与 $y\in[c/\alpha_N,C/\alpha_N]$，有
\begin{equation}\label{eq:W}
 W_N(y)=\frac12e^{2y}\Bigl(y-\frac12\Bigr)\bigl(1+o(1)\bigr),
\end{equation}
且在该范围内一致。
\end{lemma}

\begin{proof}
取 $x=e^y$，由引理~\ref{lem:pnt}，
\[
 W_N(y)=\frac12e^{2y}y-\frac14e^{2y}+O\bigl(e^{2y}y^Ce^{-c\sqrt{y}}\bigr)
 =e^{2y}\Bigl(\frac y2-\frac14\Bigr)(1+o(1)),
\]
而 $\frac y2-\frac14=\frac12(y-\frac12)$。低阶项即注~\ref{rem:primepowers}
中的项；误差在 $y\in[c/\alpha_N,C/\alpha_N]$ 上一致，因为
$e^{-c\sqrt{y}}\to0$ 而 $y^C$ 被 $\alpha_N^{-1}$ 的幂控制。
\end{proof}

\subsection*{3.3.\ 模型核}

在 \eqref{eq:W} 成立的区域中，$\sqrt{m(y)m(u)}$ 里的第二边界项指数小：
\begin{equation}\label{eq:second}
 \sqrt{m(y)m(u)}
 =e^{-(1+\alpha_N)(y+u)}
 \bigl(1+O(e^{-2(\beta-\alpha_N)\min(y,u)})\bigr).
\end{equation}
事实上 $m(y)=e^{-2(1+\alpha_N)y}(1+e^{-2(\beta-\alpha_N)y})$，而当
$y\ge c/\alpha_N$ 时 $e^{-2(\beta-\alpha_N)y}\le e^{-2(\beta-\alpha_N)c/\alpha_N}\to0$。
这引出如下定义。

\begin{definition}\label{def:model}
模型核为
\begin{equation}\label{eq:model}
 \mathscr A_N(y,u)=\pi\Bigl(\min(y,u)-\frac12\Bigr)
 e^{-2\alpha_N\min(y,u)}e^{-(1+\alpha_N)|y-u|},
 \qquad y,u\ge y_0 .
\end{equation}
\end{definition}

\begin{lemma}\label{lem:close}
对固定的 $0<c<C<\infty$ 与 $y,u\in I_0=[c/\alpha_N,C/\alpha_N]$，
\begin{equation}\label{eq:close}
 \mathscr D_N(y,u)=\mathscr A_N(y,u)\bigl(1+O(\epsilon_N(c,C))\bigr),
 \qquad \epsilon_N(c,C)\longrightarrow0 .
\end{equation}
\end{lemma}

\begin{proof}
在 $I_0$ 上，引理~\ref{lem:W} 一致地给出
$W_N(\min)=\frac12e^{2\min}(\min-\frac12)(1+o(1))$，而 \eqref{eq:second}
给出第二个因子。相乘得
\[
 \begin{aligned}
 \mathscr D_N
 &=2\pi\cdot\frac12 e^{2\min}\Bigl(\min-\frac12\Bigr)e^{-(1+\alpha_N)(y+u)}(1+o(1))\\
 &=\pi\Bigl(\min-\frac12\Bigr)e^{2\min-(1+\alpha_N)(y+u)}(1+o(1)).
 \end{aligned}
\]
由于 $2\min-(1+\alpha_N)(y+u)=-2\alpha_N\min-(1+\alpha_N)|y-u|$，这正是
\eqref{eq:model}，相对误差为 $o(1)$。
\end{proof}

\section{模型算子的范数}\label{sec:model}

本节中 $\mathscr A_N$ 始终表示 $[y_0,\infty)\times[y_0,\infty)$ 上的核
\eqref{eq:model}，并视为 $L^2([y_0,\infty))$ 上的算子。

\subsection*{4.1.\ Schur 检验}

置
\begin{equation}\label{eq:Valpha}
 V_\alpha(y)=\Bigl(y-\frac12\Bigr)e^{-2\alpha_Ny},\qquad y\ge y_0 .
\end{equation}
则 $\mathscr A_N(y,u)=\pi V_\alpha(\min(y,u))e^{-(1+\alpha_N)|y-u|}$。

\begin{lemma}\label{lem:schur}
对 $u\le y$，
\begin{equation}\label{eq:ratio}
 \frac{V_\alpha(u)}{V_\alpha(y)}e^{-(1+\alpha_N)(y-u)}
 \le e^{-(1-\alpha_N)(y-u)},
\end{equation}
从而
\begin{equation}\label{eq:schur}
 \sup_{y\ge y_0}\int_{y_0}^{\infty}\mathscr A_N(y,u)\dd u
 \le\pi\|V_\alpha\|_\infty
 \Bigl(\frac1{1-\alpha_N}+\frac1{1+\alpha_N}\Bigr).
\end{equation}
\end{lemma}

\begin{proof}
对 $u\le y$，有 $(u-\frac12)/(y-\frac12)\le1$ 与
$e^{-2\alpha_N(u-y)}=e^{2\alpha_N(y-u)}$，这给出 \eqref{eq:ratio}。于是
\[
 \int_{y_0}^{y}\mathscr A_N(y,u)\dd u
 \le\pi\|V_\alpha\|_\infty\int_{y_0}^{y}e^{-(1-\alpha_N)(y-u)}\dd u
 \le\frac{\pi\|V_\alpha\|_\infty}{1-\alpha_N},
\]
且
\[
 \int_{y}^{\infty}\mathscr A_N(y,u)\dd u
 =\pi V_\alpha(y)\int_0^{\infty}e^{-(1+\alpha_N)v}\dd v
 \le\frac{\pi\|V_\alpha\|_\infty}{1+\alpha_N}.
\]
两式相加即得 \eqref{eq:schur}。
\end{proof}

\subsection*{4.2.\ 峰值}

\begin{lemma}\label{lem:peak}
函数 $V_\alpha$ 在 $[y_0,\infty)$ 上的最大值在
\begin{equation}\label{eq:ystar}
 y_D^*=\frac12+\frac1{2\alpha_N}
\end{equation}
处取得，且
\begin{equation}\label{eq:peak}
 V_\alpha(y_D^*)=\frac{e^{-1-\alpha_N}}{2\alpha_N}.
\end{equation}
\end{lemma}

\begin{proof}
$V_\alpha'(y)=e^{-2\alpha_Ny}[1-2\alpha_N(y-\frac12)]$ 在 $y_D^*$ 处为零，
且在该点由正变负；对大 $N$ 有 $y_D^*>y_0$。代入即得数值。
\end{proof}

\begin{proposition}\label{prop:schur}
有
\begin{equation}\label{eq:upper}
 \|\mathscr A_N\|\le\frac{\pi e^{-1-\alpha_N}}{\alpha_N(1-\alpha_N^2)}
 =\frac{\pi}{e\alpha_N}(1+o(1)).
\end{equation}
\end{proposition}

\begin{proof}
把 \eqref{eq:schur} 与 \eqref{eq:peak} 结合。
\end{proof}

\subsection*{4.3.\ 匹配下界}

\begin{lemma}\label{lem:lower}
有
\begin{equation}\label{eq:lower}
 \|\mathscr A_N\|\ge\frac{\pi e^{-1-\alpha_N}}{2\alpha_N}
 \cdot\frac{2}{1+\alpha_N}(1-o(1))
 =\frac{\pi}{e\alpha_N}(1-o(1)).
\end{equation}
\end{lemma}

\begin{proof}
置 $B_N=\alpha_N^{-1/2}$ 与
\[
 J_N=\Bigl[y_D^*-\frac{B_N}2,\;y_D^*+\frac{B_N}2\Bigr],
 \qquad f_N=B_N^{-1/2}\mathbf 1_{J_N},
\]
则 $\|f_N\|_2=1$。在 $J_N$ 上，
\[
 V_\alpha(y)=V_\alpha(y_D^*)\bigl(1-O(\alpha_N^2(y-y_D^*)^2)\bigr)
 =V_\alpha(y_D^*)(1-O(\alpha_N)),
\]
因为 $\alpha_NB_N^2=1$ 且 $V_\alpha''(y_D^*)/V_\alpha(y_D^*)=-4\alpha_N^2$。
故
\[
 \begin{aligned}
 \langle\mathscr A_Nf_N,f_N\rangle
 &=\frac1{B_N}\iint_{J_N^2}\mathscr A_N(y,u)\dd y\dd u\\
 &\ge\pi V_\alpha(y_D^*)(1-O(\alpha_N))
 \cdot\frac1{B_N}\iint_{J_N^2}e^{-(1+\alpha_N)|y-u|}\dd y\dd u .
 \end{aligned}
\]
该二重积分等于
\[
 \frac2{1+\alpha_N}\Bigl(B_N-\frac{1-e^{-(1+\alpha_N)B_N}}{1+\alpha_N}\Bigr)
 =\frac{2B_N}{1+\alpha_N}(1-o(1)),
\]
故商为 $\frac{2}{1+\alpha_N}(1-o(1))$，断言由 \eqref{eq:peak} 得出。
\end{proof}

\begin{proposition}\label{prop:modelnorm}
有
\begin{equation}\label{eq:modelnorm}
 \|\mathscr A_N\|=\Bigl(\frac{\pi}{e}+o(1)\Bigr)\alpha_N^{-1}.
\end{equation}
\end{proposition}

\begin{proof}
结合命题~\ref{prop:schur} 与引理~\ref{lem:lower}。
\end{proof}

\section{比较}\label{sec:compare}

\subsection*{5.1.\ 三个区域}

固定 $0<c<\frac14$ 与 $C>1$，置
\begin{equation}\label{eq:regions}
 I_-=[y_0,c/\alpha_N],\qquad I_0=[c/\alpha_N,C/\alpha_N],\qquad
 I_+=[C/\alpha_N,\infty),
\end{equation}
其中 $y_0$ 如 \eqref{eq:y0}。设 $P_-,P_0,P_+$ 是 $L^2([y_0,\infty))$ 到这些
区间的正交投影。我们将逐块比较 $\mathcal D_N$ 与 $\mathscr A_N$。

\subsection*{5.2.\ 行估计}

\begin{lemma}\label{lem:row}
对 $y\ge y_0$，
\begin{equation}\label{eq:WNbound}
 W_N(y)\ll e^{2y}(1+y),
\end{equation}
且
\begin{equation}\label{eq:Dbound}
 \mathscr D_N(y,u)\ll(1+\min(y,u))e^{-2\alpha_N\min(y,u)}
 e^{-(1+\alpha_N)|y-u|}.
\end{equation}
\end{lemma}

\begin{proof}
对第一个界，用 $q\le e^y$、当 $n\le x$ 时 $\Lambda(n)^2\le(\log x)\Lambda(n)$
以及 $\sum_{n\le x}\Lambda(n)\ll x$：
\[
 W_N(y)\le e^y\sum_{n\le e^y}\Lambda(n)^2
 \ll e^y\cdot e^y\cdot(1+y)=e^{2y}(1+y).
\]
对第二个界，把 \eqref{eq:WNbound} 与
$\sqrt{m(y)m(u)}\le2e^{-(1+\alpha_N)(y+u)}$ 结合，并用
$2\min-(1+\alpha_N)(y+u)=-2\alpha_N\min-(1+\alpha_N)|y-u|$。
\end{proof}

\begin{lemma}\label{lem:rows}
对 $y\ge y_0$，
\begin{equation}\label{eq:rowD}
 \int_{y_0}^{\infty}\mathscr D_N(y,u)\dd u\ll R_\alpha(y),
 \qquad
 \int_{y_0}^{\infty}\mathscr A_N(y,u)\dd u\ll R_\alpha(y),
\end{equation}
其中 $R_\alpha(y)=(1+y)e^{-2\alpha_Ny}$。
\end{lemma}

\begin{proof}
对 $\mathscr D_N$，把 \eqref{eq:Dbound} 分别对 $u\le y$ 与 $u\ge y$ 积分：
\[
 \int_{y_0}^{y}\mathscr D_N\dd u\ll R_\alpha(y)\int_{y_0}^{y}
 e^{-(1-\alpha_N)(y-u)}\dd u\ll\frac{R_\alpha(y)}{1-\alpha_N},
\]
\[
 \int_{y}^{\infty}\mathscr D_N\dd u\ll R_\alpha(y)
 \int_0^{\infty}e^{-(1+\alpha_N)v}\dd v\ll\frac{R_\alpha(y)}{1+\alpha_N}.
\]
对 $\mathscr A_N$，同样的计算以 $V_\alpha\le R_\alpha$ 与比值界 \eqref{eq:ratio}
成立。
\end{proof}

\subsection*{5.3.\ 边缘块}

\begin{lemma}\label{lem:edge}
有
\begin{equation}\label{eq:low}
 \|P_-\mathcal D_NP_-\|+\|P_-\mathscr A_NP_-\|\ll1+c/\alpha_N,
\end{equation}
且
\begin{equation}\label{eq:high}
 \|P_+\mathcal D_NP_+\|+\|P_+\mathscr A_NP_+\|\ll Ce^{-2C}/\alpha_N .
\end{equation}
\end{lemma}

\begin{proof}
对低区，注意当 $y\le c/\alpha_N$ 时 $R_\alpha(y)\le1+y\le1+c/\alpha_N$；
用行界 \eqref{eq:rowD} 作 Schur 检验即得 \eqref{eq:low}。对高区，$R_\alpha$ 在
$y\ge C/\alpha_N$ 上递减（因为当 $2\alpha_N(1+y)>1$ 时
$R_\alpha'(y)=e^{-2\alpha_Ny}[1-2\alpha_N(1+y)]<0$，而 $C>1/2$ 保证了这一点），故
\[
 \sup_{y\ge C/\alpha_N}R_\alpha(y)=R_\alpha(C/\alpha_N)
 \le\Bigl(1+\frac C{\alpha_N}\Bigr)e^{-2C}\ll\frac{Ce^{-2C}}{\alpha_N},
\]
再由 Schur 检验得 \eqref{eq:high}。
\end{proof}

\subsection*{5.4.\ 中心块}

\begin{lemma}\label{lem:center}
对固定的 $c,C$，
\begin{equation}\label{eq:center}
 \|P_0(\mathcal D_N-\mathscr A_N)P_0\|\le\epsilon_N(c,C)\|P_0\mathscr A_NP_0\|
 =o(\alpha_N^{-1}),
\end{equation}
其中 $\epsilon_N(c,C)\to0$。
\end{lemma}

\begin{proof}
由引理~\ref{lem:close}，在 $I_0\times I_0$ 上
\[
 |\mathscr D_N-\mathscr A_N|\le\epsilon_N(c,C)\mathscr A_N .
\]
由于 $\mathscr A_N$ 在 $I_0$ 上 $\ge0$（见下文引理~\ref{lem:positive}），对任意
$f\in L^2(I_0)$，记 $|f|$ 为 $f$ 的模，则
\[
 \bigl|\langle(\mathcal D_N-\mathscr A_N)f,f\rangle\bigr|
 \le\epsilon_N\int_{I_0}\!\int_{I_0}\mathscr A_N(y,u)\,|f(y)|\,|f(u)|\dd y\dd u
 =\epsilon_N\langle\mathscr A_N|f|,|f|\rangle
 \le\epsilon_N\|\mathscr A_N\|\,\|f\|^2 .
\]
对 $\|f\|=1$ 取上确界（注意 $f\mapsto|f|$ 保持 $L^2$ 范数），即得算子不等式
$\|P_0(\mathcal D_N-\mathscr A_N)P_0\|\le\epsilon_N\|P_0\mathscr A_NP_0\|$。
最后，由 \eqref{eq:rowD} 有 $\|P_0\mathscr A_NP_0\|\ll\alpha_N^{-1}$。
\end{proof}

\subsection*{5.5.\ 模型的正性}

\begin{lemma}\label{lem:positive}
核 $\mathscr A_N$ 在 $[y_0,\infty)^2$ 上半正定。
\end{lemma}

\begin{proof}
写
\[
 \mathscr A_N(y,u)=\pi e^{-(1+\alpha_N)(y+u)}U(\min(y,u)),
 \qquad U(y)=\Bigl(y-\frac12\Bigr)e^{2y}.
\]
由于 $y_0>\frac12$，有 $U(y_0)>0$，且当 $y\ge y_0$ 时 $U'(y)=2ye^{2y}>0$。故
\[
 U(\min(y,u))=U(y_0)+\int_{y_0}^{\infty}U'(v)\mathbf 1_{y\ge v}\mathbf 1_{u\ge v}\dd v .
\]
第一项给出秩一核 $U(y_0)e^{-(1+\alpha_N)(y+u)}$，第二项给出秩一核
$e^{-(1+\alpha_N)(y+u)}\mathbf 1_{y\ge v}\mathbf 1_{u\ge v}$ 的非负组合。故
$\mathscr A_N$ 是正核。
\end{proof}

\subsection*{5.6.\ 交叉块}

\begin{lemma}\label{lem:cross}
设 $T\ge0$ 是希尔伯特空间上的有界算子，$P,Q$ 为正交投影。则
\begin{equation}\label{eq:cross}
 \|PTQ\|^2\le\|PTP\|\,\|QTQ\| .
\end{equation}
\end{lemma}

\begin{proof}
写 $PTQ=(P T^{1/2})(T^{1/2}Q)$，用 $\|AB\|\le\|A\|\|B\|$ 以及
$\|PT^{1/2}\|^2=\|PTP\|$、$\|T^{1/2}Q\|^2=\|QTQ\|$。
\end{proof}

\subsection*{5.7.\ 比较}

\begin{proposition}\label{prop:compare}
对固定的 $c,C$，
\begin{equation}\label{eq:compare}
 \limsup_{N\to\infty}\alpha_N\|\mathcal D_N-\mathscr A_N\|
 \ll c+\sqrt c+Ce^{-2C}+\sqrt{Ce^{-2C}}+\sqrt{cCe^{-2C}} .
\end{equation}
\end{proposition}

\begin{proof}
由引理~\ref{lem:positive}，$\mathscr A_N\ge0$；由引理~\ref{lem:kernel}，
$\mathcal D_N\ge0$。对 $T=\mathcal D_N$ 与 $T=\mathscr A_N$，引理~\ref{lem:cross}
与对角估计 \eqref{eq:low}--\eqref{eq:center} 给出，对 $i\ne j$，
\[
 \|P_iTP_j\|\ll\frac{\sqrt{c}}{\alpha_N},\quad
 \|P_0TP_+\|\ll\frac{\sqrt{Ce^{-2C}}}{\alpha_N},\quad
 \|P_-TP_+\|\ll\frac{\sqrt{cCe^{-2C}}}{\alpha_N} .
\]
$\mathcal D_N-\mathscr A_N$ 的九个块被 $\mathcal D_N$ 与 $\mathscr A_N$ 的对应块
之和控制（三角不等式），故
\[
 \alpha_N\|\mathcal D_N-\mathscr A_N\|
 \ll c+\sqrt c+Ce^{-2C}+\sqrt{Ce^{-2C}}+\sqrt{cCe^{-2C}}+\epsilon_N(c,C).
\]
对 $N$ 取 $\limsup$ 并用 $\epsilon_N(c,C)\to0$ 即得 \eqref{eq:compare}。
\end{proof}

\section{主定理的证明}\label{sec:proof}

\begin{proof}[定理~\ref{thm:main} 的证明]
取 $0<c<\frac14$、$C>1$。由命题~\ref{prop:compare}，
\[
 \limsup_{N\to\infty}\alpha_N\|\mathcal D_N-\mathscr A_N\|\ll\delta(c,C),
\]
其中 $\delta(c,C)=c+\sqrt c+Ce^{-2C}+\sqrt{Ce^{-2C}}+\sqrt{cCe^{-2C}}$。
由于当 $c\downarrow0$、$C\uparrow\infty$ 时 $\delta(c,C)\to0$，我们得到
\[
 \alpha_N\|\mathcal D_N-\mathscr A_N\|\longrightarrow0 .
\]
由命题~\ref{prop:modelnorm}，
$\|\mathscr A_N\|=(\pi/e+o(1))\alpha_N^{-1}$，因此
\[
 \|\mathcal D_N\|=\|\mathscr A_N\|+O\bigl(\|\mathcal D_N-\mathscr A_N\|\bigr)
 =\Bigl(\frac{\pi}{e}+o(1)\Bigr)\alpha_N^{-1}
 =\Bigl(\frac{\pi}{e}+o(1)\Bigr)\log N .
\]
\end{proof}

\begin{proof}[推论~\ref{cor:HS} 的证明]
由于 $\mathcal D_N\ge0$ 且 $\operatorname{tr}\mathcal D_N\ll\alpha_N^{-2}$，
\[
 \|\mathcal D_N\|_{\HS}^2=\operatorname{tr}(\mathcal D_N^2)
 \le\|\mathcal D_N\|\operatorname{tr}\mathcal D_N
 \ll\alpha_N^{-1}\cdot\alpha_N^{-2}=\alpha_N^{-3},
\]
故 $\|\mathcal D_N\|_{\HS}\ll\alpha_N^{-3/2}=(\log N)^{3/2}$。
\end{proof}

\section{注}\label{sec:remarks}

\begin{remark}\label{rem:constant}
常数 $\pi/e$ 有明确的来源。$\mathscr D_N$ 在对角线上的局部极限是以
$e^{-(1+\alpha_N)|y-u|}$ 为核的卷积，其符号
$2(1+\alpha_N)/((1+\alpha_N)^2+t^2)$ 在 $t=0$ 处取最大值 $2/(1+\alpha_N)$；
乘子 $\pi$ 来自权重，指数 $e^{-1-\alpha_N}/(2\alpha_N)$ 来自峰值 $y_D^*$。
乘积 $2\pi\cdot\frac12\cdot e^{-1}$ 即 $\pi/e$。
\end{remark}

\begin{remark}\label{rem:vsLL}
论文 \cite[命题 6.5]{LL} 只用到
$\|\mathcal D_N\|_{\HS}\le\operatorname{tr}\mathcal D_N\ll(\log N)^2$。
定理~\ref{thm:main} 给出带精确常数的算子范数，推论~\ref{cor:HS} 把
希尔伯特--施密特界改进到 $(\log N)^{3/2}$。
\end{remark}

\begin{remark}\label{rem:limitation}
论证只用到素数定理；引理~\ref{lem:pnt} 的误差项除一致性外并不需要。特别地，
结论是无条件的。它只涉及 $\mathcal D_N$ 自身的范数，不使用、也不声称关于
\eqref{eq:split} 中完整算子 $T_N$ 的谱的任何信息。
\end{remark}

\begin{thebibliography}{9}

\bibitem[LL]{LL}
J. Li and J. Liu,
\emph{A conditional proof of the Goldbach conjecture},
arXiv:2606.05224 (2026).

\end{thebibliography}

\end{document}
